\section{Reproducibility and Evidence Access}
The accompanying package stores experiment, controller, physics-configuration, claim, figure-source, and comparability registries. Each primary numerical claim has an identifier linking its manuscript location to an immutable artifact and a stated calculation. Missing source fields remain empty. Original results are referenced in place and verified by SHA-256; checkpoint files are not edited or redistributed as newly trained models.

Offline figure generation reads archived NPZ/JSON/CSV files. It performs no native time integration. The parameter extractor compiles an archived MJCF into a model for reading inertial arrays but does not construct a new rollout. Existing regression tests run separately under bounded timeouts. Commands and software metadata are recorded in the reproducibility directory. The historical model, contact adapter, reward, seed registries, and training protocols remain unchanged.

Full-horizon figures use L seed1801001 and its registered cohort; matched learning plots use the Y final policies and development contexts. U nested windows use saved native prefixes, and their dependence is explicit. Historical J/K records remain visible without being pooled into the later contact domain. Readers should reproduce plots in a fresh output namespace, preserve original data hashes, and distinguish data-processing reproducibility from rerunning expensive experiments.

LaTeX source uses a generic two-column article layout rather than claiming compliance with a named venue. Author identities, affiliations, data-release permissions, and target publication requirements must be supplied or approved by the project owner. Compilation availability and unresolved evidence gaps are reported in the handoff; a compilable document alone would not establish submission readiness.
